\documentclass[10pt,a4paper]{amsart}
\usepackage{amsmath,amssymb,amsthm,graphicx}
% For commutative diagrams we use Xy-pic
\usepackage[all]{xy}

% We use 2cell for 2-commutative diagrams.
\xyoption{2cell}
% Theorem environments.
%
\theoremstyle{plain}
\newtheorem{theorem}[subsection]{Theorem}
\newtheorem{proposition}[subsection]{Proposition}
\newtheorem{lemma}[subsection]{Lemma}

\theoremstyle{definition}
\newtheorem{definition}[subsection]{Definition}
\newtheorem{example}[subsection]{Example}
\newtheorem{exercise}[subsection]{Exercise}
\newtheorem{situation}[subsection]{Situation}

\theoremstyle{remark}
\newtheorem{remark}[subsection]{Remark}
\newtheorem{remarks}[subsection]{Remarks}

\numberwithin{equation}{subsection}

% Macros
%
\def\lim{\mathop{\mathrm{lim}}\nolimits}
\def\colim{\mathop{\mathrm{colim}}\nolimits}
\def\Spec{\mathop{\mathrm{Spec}}}
\def\Hom{\mathop{\mathrm{Hom}}\nolimits}
\def\Ext{\mathop{\mathrm{Ext}}\nolimits}
\def\SheafHom{\mathop{\mathcal{H}\!\mathit{om}}\nolimits}
\def\SheafExt{\mathop{\mathcal{E}\!\mathit{xt}}\nolimits}
\def\Sch{\mathit{Sch}}
\def\Mor{\mathop{\mathrm{Mor}}\nolimits}
\def\Ob{\mathop{\mathrm{Ob}}\nolimits}
\def\Sh{\mathop{\mathit{Sh}}\nolimits}
\def\NL{\mathop{N\!L}\nolimits}
\def\CH{\mathop{\mathrm{CH}}\nolimits}
\def\proetale{{pro\text{-}\acute{e}tale}}
\def\etale{{\acute{e}tale}}
\def\QCoh{\mathit{QCoh}}
\def\Ker{\mathop{\mathrm{Ker}}}
\def\Im{\mathop{\mathrm{Im}}}
\def\Coker{\mathop{\mathrm{Coker}}}
\def\Coim{\mathop{\mathrm{Coim}}}

% Boxtimes
%
\DeclareMathSymbol{\boxtimes}{\mathbin}{AMSa}{"02}

%
% Macros for moduli stacks/spaces
%
\def\QCohstack{\mathcal{QC}\!\mathit{oh}}
\def\Cohstack{\mathcal{C}\!\mathit{oh}}
\def\Spacesstack{\mathcal{S}\!\mathit{paces}}
\def\Quotfunctor{\mathrm{Quot}}
\def\Hilbfunctor{\mathrm{Hilb}}
\def\Curvesstack{\mathcal{C}\!\mathit{urves}}
\def\Polarizedstack{\mathcal{P}\!\mathit{olarized}}
\def\Complexesstack{\mathcal{C}\!\mathit{omplexes}}
% \Pic is the operator that assigns to X its picard group, usage \Pic(X)
% \Picardstack_{X/B} denotes the Picard stack of X over B
% \Picardfunctor_{X/B} denotes the Picard functor of X over B
\def\Pic{\mathop{\mathrm{Pic}}\nolimits}
\def\Picardstack{\mathcal{P}\!\mathit{ic}}
\def\Picardfunctor{\mathrm{Pic}}
\def\Deformationcategory{\mathcal{D}\!\mathit{ef}}
\usepackage{cleveref}
\crefname{originallemma}{original lemma}{original lemmas}
\Crefname{originallemma}{Original lemma}{Original lemmas}
\makeatletter
\newlabel{lemma-finite-periodic-length}{{0EA8}{source}{}{sourcepointer-1}{}}
\newlabel{lemma-finite-periodic-length@cref}{{[originallemma][0][]0EA8}{[1][0][]source}{}{}{}}
\makeatother
\crefname{originallemma}{original lemma}{original lemmas}
\Crefname{originallemma}{Original lemma}{Original lemmas}
\makeatletter
\newlabel{algebra-lemma-length-additive}{{00IV}{source}{}{sourcepointer-2}{}}
\newlabel{algebra-lemma-length-additive@cref}{{[originallemma][0][]00IV}{[1][0][]source}{}{}{}}
\makeatother
\crefname{originallemma}{original lemma}{original lemmas}
\Crefname{originallemma}{Original lemma}{Original lemmas}
\makeatletter
\newlabel{lemma-additivity-periodic-length}{{0EA7}{source}{}{sourcepointer-3}{}}
\newlabel{lemma-additivity-periodic-length@cref}{{[originallemma][0][]0EA7}{[1][0][]source}{}{}{}}
\makeatother
\crefname{originallemma}{original lemma}{original lemmas}
\Crefname{originallemma}{Original lemma}{Original lemmas}
\makeatletter
\newlabel{algebra-lemma-Noetherian-power-ideal-kills-module}{{00L6}{source}{}{sourcepointer-4}{}}
\newlabel{algebra-lemma-Noetherian-power-ideal-kills-module@cref}{{[originallemma][0][]00L6}{[1][0][]source}{}{}{}}
\makeatother
\crefname{originaldefinition}{original definition}{original definitions}
\Crefname{originaldefinition}{Original definition}{Original definitions}
\makeatletter
\newlabel{algebra-definition-lattice}{{02ME}{source}{}{sourcepointer-5}{}}
\newlabel{algebra-definition-lattice@cref}{{[originaldefinition][0][]02ME}{[1][0][]source}{}{}{}}
\makeatother
\crefname{originaldefinition}{original definition}{original definitions}
\Crefname{originaldefinition}{Original definition}{Original definitions}
\makeatletter
\newlabel{algebra-definition-distance}{{02MG}{source}{}{sourcepointer-6}{}}
\newlabel{algebra-definition-distance@cref}{{[originaldefinition][0][]02MG}{[1][0][]source}{}{}{}}
\makeatother
\crefname{originallemma}{original lemma}{original lemmas}
\Crefname{originallemma}{Original lemma}{Original lemmas}
\makeatletter
\newlabel{algebra-lemma-order-vanishing-determinant}{{02MI}{source}{}{sourcepointer-7}{}}
\newlabel{algebra-lemma-order-vanishing-determinant@cref}{{[originallemma][0][]02MI}{[1][0][]source}{}{}{}}
\makeatother
\title{The multiplicity of a one-dimensional module after multiplying by an element, as a sum over the height one primes}
\author{The Stacks Project Authors}
\date{stacks-project:chow.tex:local-snapshot-20261012; GNU FDL 1.2 or later}
\begin{document}
\maketitle
\noindent\textit{Editorial note.} All delivered mathematical text is the author's own, copied from the pinned Stacks Project chapter file chow.tex; no statement or proof step was rewritten, reordered or removed. The author's typos 'lemma for 2 out of 3' and 'additivty' are reproduced verbatim.\par\medskip
\noindent\textit{Editorial note.} source\_id and source\_version carry the current authoritative snapshot label of this source, 20261012, read from sources/stacks-chow/license-observation.json and SOURCE.json at the moment of opening (the source was re-stamped during the night; earlier chow cards keep 20261011 and are recorded in agent-authoring/stk/license-known-deltas.md).\par\medskip
\noindent\textit{Editorial note.} The preamble is the chapter's own native preamble excerpt used for this source: lines 14-18 (Xy-pic setup) and 154-217 (theorem declarations and mathematical macros) of sources/stacks-chow/preamble.tex. The quoted class line and the xr-hyper externaldocument block are not included because the assembler supplies a plain amsart wrapper; no macro was invented or redefined.\par\medskip
\noindent\textit{Editorial note.} Same-chapter references are located at their \textbackslash{}label lines inside the pinned chapter file; the cross-chapter pointers are located in the project's own tag list staged in the snapshot (sources/stacks-chow/tags.txt), at the verbatim lines carrying their canonical tags.\par\medskip
\noindent\textit{Allowed standard prerequisites (editorial).}\begin{itemize}
\item The multiplicity of a periodic complex of finite length modules, and additivity in short exact sequences. Same chapter, Lemmas lemma-finite-periodic-length (source line 258, Stacks tag 0EA8) and lemma-additivity-periodic-length (source line 209, Stacks tag 0EA7); invoked at proof lines 332 and 337. Supply the vanishing of both sides in the finite length case and the additivity of multiplicities used twice to cut the problem down to the one-dimensional case.
\item Additivity of the length function. Algebra, Lemma labelled algebra-lemma-length-additive (canonical Stacks tag 00IV); invoked at proof lines 335-336. Supplies the other half of the two-out-of-three argument and later the additivity along the filtration.
\item A finite module over a Noetherian ring is killed by a power of an ideal contained in its annihilator, lattices and their distance. Algebra, Lemma algebra-lemma-Noetherian-power-ideal-kills-module (00L6, line 349), Definition algebra-definition-lattice (02ME, line 360) and Definition algebra-definition-distance (02MG, line 363). Supply the q-adic filtration reducing to q M = 0, and the dictionary between length\_R(M/xM) and the distance d(M, xM) for a lattice M inside a one-dimensional fraction field.
\item The order of vanishing of a determinant computes the length of the cokernel. Algebra, Lemma labelled algebra-lemma-order-vanishing-determinant (canonical Stacks tag 02MI); invoked at proof lines 367-368. Supplies the final equality in the one-dimensional lattice case, where x acts injectively and the determinant is x\textasciicircum{}\{dim\_K(V)\}.
\end{itemize}\medskip
\ifdefined\hypertarget\hypertarget{sourcepointer-1}{}\fi
\noindent\textit{Source reference.} Original lemma 0EA8: sources/stacks-chow/pinned-original-source.tex, source lines 258--258. Same-chapter pointer at proof line 332 for the finite length case; the reduction argument built on it is retained in the counted span.\par\medskip
\ifdefined\hypertarget\hypertarget{sourcepointer-2}{}\fi
\noindent\textit{Source reference.} Original lemma 00IV: sources/stacks-chow/tags.txt, source lines 693--693. Cross-chapter standard prerequisite invoked at proof lines 335-336 for additivity of length.\par\medskip
\ifdefined\hypertarget\hypertarget{sourcepointer-3}{}\fi
\noindent\textit{Source reference.} Original lemma 0EA7: sources/stacks-chow/pinned-original-source.tex, source lines 209--209. Same-chapter pointer at proof line 337 for additivity of multiplicities, the second half of the two-out-of-three step.\par\medskip
\ifdefined\hypertarget\hypertarget{sourcepointer-4}{}\fi
\noindent\textit{Source reference.} Original lemma 00L6: sources/stacks-chow/tags.txt, source lines 776--776. Cross-chapter standard prerequisite invoked once at proof line 349 for the termination of the q-adic filtration.\par\medskip
\ifdefined\hypertarget\hypertarget{sourcepointer-5}{}\fi
\noindent\textit{Source reference.} Original definition 02ME: sources/stacks-chow/tags.txt, source lines 3441--3441. Cross-chapter pointer at proof line 360 for the notion of a lattice used to cut by torsion.\par\medskip
\ifdefined\hypertarget\hypertarget{sourcepointer-6}{}\fi
\noindent\textit{Source reference.} Original definition 02MG: sources/stacks-chow/tags.txt, source lines 3443--3443. Cross-chapter pointer at proof line 363 for the distance d(M, xM), the translation step of the argument.\par\medskip
\ifdefined\hypertarget\hypertarget{sourcepointer-7}{}\fi
\noindent\textit{Source reference.} Original lemma 02MI: sources/stacks-chow/tags.txt, source lines 3445--3445. Cross-chapter standard prerequisite invoked once at proof lines 367-368 for the final determinant computation.\par\medskip
\begin{lemma}
\label{lemma-length-multiplication}
Let $R$ be a Noetherian local ring.
Let $M$ be a finite $R$-module. Let $x \in R$. Assume that
\begin{enumerate}
\item $\dim(\text{Supp}(M)) \leq 1$, and
\item $\dim(\text{Supp}(M/xM)) \leq 0$.
\end{enumerate}
Write
$\text{Supp}(M) = \{\mathfrak m, \mathfrak q_1, \ldots, \mathfrak q_t\}$.
Then
$$
e_R(M, 0, x) =
\sum\nolimits_{i = 1, \ldots, t}
\text{ord}_{R/\mathfrak q_i}(x)
\text{length}_{R_{\mathfrak q_i}}(M_{\mathfrak q_i}).
$$
\end{lemma}

\begin{proof}
We first make some preparatory remarks.
The result of the lemma holds if $M$ has finite length, i.e., if $t = 0$,
because both the left hand side and the right hand side are zero
in this case, see Lemma \ref{lemma-finite-periodic-length}.
Also, if we have a short exact sequence $0 \to M \to M' \to M'' \to 0$
of modules satisfying (1) and (2), then lemma for 2 out of 3
of these implies the lemma for the third by the
additivity of length (Algebra, Lemma \ref{algebra-lemma-length-additive}) and
additivty of multiplicities (Lemma \ref{lemma-additivity-periodic-length}).

\medskip\noindent
Denote $M_i$ the image of $M$ in $M_{\mathfrak q_i}$, so
$\text{Supp}(M_i) = \{\mathfrak m, \mathfrak q_i\}$.
The kernel and cokernel of the map $M \to \bigoplus M_i$
have support $\{\mathfrak m\}$ and hence have finite length.
By our preparatory remarks, it follows that it suffices to
prove the lemma for each $M_i$. Thus we may assume that
$\text{Supp}(M) = \{\mathfrak m, \mathfrak q\}$.
In this case we have a finite filtration
$M \supset \mathfrak qM \supset \mathfrak q^2M \supset \ldots \supset
\mathfrak q^nM = 0$ by Algebra, Lemma
\ref{algebra-lemma-Noetherian-power-ideal-kills-module}.
Again additivity shows that it suffices to prove the lemma
in the case $M$ is annihilated by $\mathfrak q$.
In this case we can view $M$ as a $R/\mathfrak q$-module,
i.e., we may assume that $R$ is a Noetherian local domain
of dimension $1$ with fraction field $K$.
Dividing by the torsion submodule, i.e., by the
kernel of $M \to M \otimes_R K = V$ (the torsion has
finite length hence is handled by our preliminary remarks)
we may assume that $M \subset V$ is a lattice
(Algebra, Definition \ref{algebra-definition-lattice}).
Then $x : M \to M$ is injective and
$\text{length}_R(M/xM) = d(M, xM)$
(Algebra, Definition \ref{algebra-definition-distance}). Since
$\text{length}_K(V) = \dim_K(V)$
we see that $\det(x : V \to V) = x^{\dim_K(V)}$ and
$\text{ord}_R(\det(x : V \to V)) = \dim_K(V) \text{ord}_R(x)$.
Thus the desired equality follows from
Algebra, Lemma \ref{algebra-lemma-order-vanishing-determinant}
in this case.
\end{proof}

\end{document}
